\documentclass[11pt]{article}

\usepackage[a4paper,margin=2.4cm]{geometry}
\usepackage{amsmath,amssymb}
\usepackage{siunitx}
\usepackage{xcolor}
\usepackage[hidelinks]{hyperref}
\usepackage{enumitem}

\title{\texttt{PTCryspMC.jl}: a fast photon-only Monte Carlo\\
\large Project overview --- what it does and why}
\author{}
\date{\today}

\begin{document}
\maketitle

\begin{abstract}
\noindent
\texttt{PTCryspMC.jl} simulates how a PET scanner records the positron activity left by a
proton field. It reads a frozen \emph{scenario} (the annihilation points and per-isotope
decay budget produced upstream) and, for a given detector, writes the list of coincidences
the scanner would measure. It is a fast, photon-only Monte Carlo: it transports the two
\SI{511}{keV} photons of each annihilation through an analytic phantom and a segmented
crystal ring, never the protons that created them. This document is a high-level companion
to the method guide \texttt{docs/pet\_simulation.tex}: it describes \emph{what} the project
does and \emph{why} each modelling choice was made, rather than re-deriving the detailed
method. The numbers it quotes are read from the code and data tables; where a parameter is
a placeholder this is stated.
\end{abstract}

\tableofcontents

% =====================================================================
\section{Motivation and scope}
\label{sec:motivation}

\paragraph{The physics question.}
A therapeutic proton beam stops inside the patient at a depth set by its energy --- the
Bragg peak --- and along its track it activates the tissue, creating short-lived
positron emitters ($^{15}$O, $^{11}$C, $^{13}$N, \dots). Those positrons annihilate and
emit pairs of \SI{511}{keV} photons, so a PET scanner placed around the patient sees a
faint image of where the beam deposited dose. Comparing that image with the planned range
is a way to verify, \emph{in vivo}, that the protons stopped where intended. The precision
of that verification depends on the detector: its spatial, energy and timing resolution all
shape the coincidence list a scanner can record.

\paragraph{What this repository does.}
\texttt{PTCryspMC.jl} answers one question only: \emph{given a fixed activity distribution
and a candidate detector, what list of coincidences does the scanner record?} It reads a
scenario (Sec.~\ref{sec:scenario}) and produces, per detector, a list-mode coincidence
(line-of-response, LOR) file --- two \mbox{3-D} hit positions, two energies, two timestamps
and a truth label per event. The simulation is run once per detector, and every detector
reads the \emph{identical} scenario, so any difference in the output comes from the detector
alone. This makes it a controlled bench for comparing scanners.

\paragraph{What it does not do.}
It does not transport protons, and it does not compute the activation: the annihilation
points and their per-isotope counts are inputs, frozen upstream
(\texttt{ptcryspg4}/\texttt{ptcrysp-scenarios}). The positron range is already folded into
the annihilation points, so the emission points need no further treatment. The downstream
analysis that consumes the coincidence list --- range precision, detector comparison, image
reconstruction --- is separate, deferred, and lives elsewhere. The repository is, in short,
the \emph{measurement} stage and nothing else.

\paragraph{Why a fast, photon-only Monte Carlo.}
A \SI{1}{Gy} in-room field produces of order $10^{8}$ decays. A full
Geant4-style simulation tracking every secondary electron and every scintillation photon
would be far too slow to run that many decays for each of several candidate scanners. The
key realisation is that, for the PET measurement, almost all of the detector response is
captured by following only the two annihilation photons: the electrons they liberate travel
sub-millimetre at \SI{511}{keV} and so deposit their energy essentially at the point of
interaction, and the scintillation light --- while it sets the timing and energy resolution
--- can be folded in as an analytic response rather than tracked photon by photon. Dropping
electron transport and optical tracking is what makes $10^{8}$ decays per detector
affordable, and it is the central design choice of the whole project.

% =====================================================================
\section{Input: a frozen scenario}
\label{sec:scenario}

A scenario is a versioned, append-only directory read from \texttt{ptcrysp-scenarios}. It
carries:
\begin{itemize}[nosep]
\item \texttt{emitters.csv} --- the annihilation points and the isotope of each emitter:
the \emph{spatial} source.
\item \texttt{sampling\_budget\_<scenario>.csv} --- the measured number of decays $N_j$ per
isotope: the \emph{normalisation}.
\item \texttt{run\_meta.csv}, \texttt{isotopes.csv}, \texttt{SCHEMA.md} --- the dose/timing
normalisation, the isotope codes, and the column meanings.
\end{itemize}
The photon set is built by drawing $N_j$ annihilation points per isotope from the scenario
\emph{with replacement}: the points supply the spatial distribution of the activity, and
$N_j$ supplies the count. Because scenarios are append-only, the scenario name is its
version, and it is stamped into every output file as provenance.

% =====================================================================
\section{The physics modelled}
\label{sec:physics}

\subsection{Annihilation: the back-to-back pair}
Each drawn point emits two \SI{511}{keV} photons in a random (isotropic) direction. They
are \emph{not} exactly anti-parallel: the residual momentum of the electron--positron pair
gives a small \emph{acollinearity}. The code (\texttt{src/source.jl}) draws the first
direction isotropically and tilts the partner of $-\mathbf{d}_1$ by a small \mbox{2-D}
Gaussian in its transverse plane, with each projected angle having a full width at half
maximum (FWHM) of about $0.5^\circ$:
\begin{equation}
\sigma_\theta = \frac{\mathrm{FWHM}}{2\sqrt{2\ln 2}}, \qquad
\mathrm{FWHM} \approx 0.5^\circ .
\end{equation}
This blurs the line of response by a small, detector-independent amount and is the
irreducible floor on PET spatial resolution from the source itself.

\subsection{Photon interactions}
A photon is transported until it is absorbed or leaves the setup. The distance to the next
interaction is drawn from the material's total macroscopic cross section, and at the
interaction point the process is chosen in proportion to the partial cross sections
(\texttt{src/sampling.jl}). At \SI{511}{keV} two processes matter:
\begin{itemize}[nosep]
\item \textbf{Compton scattering} --- the photon transfers energy to an electron and
continues with reduced energy in a new direction, both drawn from the Klein--Nishina
distribution (sampled by the Butcher--Messel composition/rejection method, as in Geant4).
The electron recoil $E-E'$ is deposited at the interaction point.
\item \textbf{Photoelectric absorption} --- the photon is absorbed and its full energy
deposited at the interaction point.
\end{itemize}
Pair production is carried in the cross-section machinery but is \emph{closed at}
\SI{511}{keV} (below the $2m_ec^2 = \SI{1.022}{MeV}$ threshold), so it never fires for the
annihilation photons. Cross sections come from NIST XCOM mass-attenuation tables, tabulated
per compound over \SIrange{10}{10000}{keV} and interpolated log--log
(\texttt{src/nist\_data.jl}, \texttt{src/materials.jl}). Because the table is per-compound,
the macroscopic cross section is simply $\Sigma = (\mu/\rho)\,\rho$ [\si{cm^{-1}}], with no
per-element mixing bookkeeping.

\subsection{The phantom and the patient-scatter background}
The patient is treated as a cylinder of \emph{water} (\SI{8}{cm} radius and half-length in
the default geometry). Water is a good stand-in for soft tissue at \SI{511}{keV}, and its
Compton scattering \emph{is} the patient-scatter background: a fraction of the photons
Compton-scatter in the phantom before reaching the ring, degrading their direction and
energy. A coincidence in which one or both photons scattered in the phantom is recorded but
labelled \texttt{scatter}, so the background is preserved in the output rather than removed.
Treating the phantom as a simple analytic solid (rather than a voxelised body) is
deliberate: the scenario already fixes \emph{where} the activity is, and a homogeneous water
cylinder reproduces the dominant attenuation and scatter physics at a fraction of the cost.

% =====================================================================
\section{Geometry}
\label{sec:geometry}

The geometry is a Geant4-style hierarchy (\texttt{src/geometry.jl}): a pure \texttt{Solid}
(shape in its own local frame), a \texttt{LogicalVolume} (solid $+$ material) and a
\texttt{PhysicalVolume} (a logical volume placed in the world frame). The available solids
are \texttt{Cylinder}, \texttt{Box}, \texttt{Sphere} and \texttt{CylShell} (a hollow
cylinder); each implements the ray interface \texttt{is\_inside},
\texttt{distance\_to\_entry}, \texttt{distance\_to\_exit} in closed form, allocation-free
on the transport hot path. The world is a large non-interacting cylinder of air enclosing
the phantom (a placed \texttt{Cylinder}) and the detector ring.

\subsection{The crystal ring as a structured block/wheel grid}
The detector is a ring of crystal, modelled as a \emph{single} cylindrical shell
(\texttt{CylShell}) rather than a thousand separate volumes. Its readout segmentation ---
the physical blocks a real scanner reads out --- is imposed as a \emph{structured grid} on
the shell: $n_\varphi$ azimuthal sectors (``blocks'') $\times$ $n_z$ axial sectors
(``wheels''), each cell spanning the full wall depth (\texttt{Scanner} in
\texttt{src/geometry.jl}). The default ring is $n_\varphi=48$, $n_z=20$ (\num{960} blocks),
inner radius \SI{38.7}{cm}, wall thickness \SI{3.7}{cm} (the radial depth = the depth of
interaction, DOI), half-length \SI{51.2}{cm}.

The block a point belongs to is fixed by simple arithmetic on its $\varphi$ and $z$,
\begin{equation}
i_\varphi = \left\lfloor \frac{\operatorname{mod}(\operatorname{atan2}(y,x),\,2\pi)}{2\pi/n_\varphi}\right\rfloor,
\qquad
i_z = \left\lfloor \frac{z+H}{2H/n_z}\right\rfloor ,
\end{equation}
and the distance to the next boundary is the nearest of the inner/outer radius, the next
$\varphi$ plane and the next $z$ plane --- all closed-form.

\paragraph{Why a structured grid.}
Stepping a photon through the segmented ring costs the same as through a single shell,
\emph{whatever the number of blocks}: there is no list of volumes to test, no acceleration
structure, just arithmetic. The block index is $O(1)$ and the boundary distances are
analytic. This is what lets the geometry scale to fine segmentation without a transport
penalty, and it is the reason the ring is one shell with a grid rather than many solids. The
small dead gaps between physical blocks are skipped for now (the crystal is treated as
continuous and each deposit is tagged with its block); they can be re-introduced later as
thin non-interacting layers once the plane-crossing distances are added.

% =====================================================================
\section{The fast Monte Carlo}
\label{sec:mc}

\subsection{One annihilation at a time}
The measured rate is modest and spread thin. For a \SI{1}{Gy} field the $\sim 10^{8}$
decays average $\sim 10^{5}/$s over the acquisition, front-loaded (peaking near
$7\times10^{5}/$s and falling with $^{15}$O's \SI{122}{s} half-life). Spread over the
$\sim 10^{3}$ crystal blocks, each block sees $\sim 10^{2}/$s, so even with a slow
($\sim\SI{1}{\micro s}$) crystal the chance of two signals overlapping in one block is
$\sim 10^{-4}$. Pileup is therefore negligible, and the simulation processes annihilations
\emph{independently}, one photon pair at a time. (The one exception --- accidental
coincidences across different decays --- cannot occur within a single event and is added in
a separate randoms pass, Sec.~\ref{sec:timing}.)

\subsection{Transport across volumes}
Each photon is carried across the volumes by the navigator (\texttt{src/navigator.jl}):
phantom (water) $\to$ air gap $\to$ ring (crystal), switching the material at each boundary.
Inside an absorbing volume the photon is advanced by the smaller of the distance to the next
boundary and the free path sampled from the cross section
(\texttt{propagate\_photon}, \texttt{src/transport.jl}); the non-interacting air mother is
never stepped, only skipped to the next interface. The production path uses an
allocation-free single-photon navigator for throughput on $10^{8}$-decay runs. Every
recorded interaction is tagged with its volume and, in the ring, its readout block.

\subsection{In the crystal: overspill and containment}
Everything that can happen in the crystal comes out of following the photon, with no special
cases:
\begin{itemize}[nosep]
\item it deposits its full energy in one block (full-energy event);
\item it Compton-scatters and the scattered photon crosses into a \emph{neighbouring} block
before being absorbed --- an \emph{overspill}, leaving energy in two blocks;
\item the scattered photon escapes the crystal altogether (out the back, the axial end, or
back toward the patient) --- only part of the \SI{511}{keV} is recorded (non-containment);
\item it Compton-scatters more than once within the same block.
\end{itemize}
The transport produces all of these by construction; the bookkeeping below turns the
deposits into hits.

\subsection{Hits: the first-interaction LOR point}
The deposits in one block are reduced to a hit (\texttt{src/coincidences.jl}): the energy is
their \emph{sum}, the time is the \emph{earliest}, and the position is the \emph{first}
interaction point --- not the energy-weighted centroid. The first interaction is the point
the line of response should pass through, and the quantity a reconstruction aims to recover
(in a monolithic crystal a CNN can separate one- from two-Compton events to find it). The
full per-interaction truth is kept in the singles record so that one- versus two-Compton
separation can be studied downstream.

% =====================================================================
\section{Detector response}
\label{sec:response}

The hit is then blurred by the detector resolution (\texttt{src/detector.jl}), and this is
where candidate scanners differ.

\paragraph{Position.}
Each hit coordinate (including the DOI) is smeared by a Gaussian of width
$\sigma_{xyz}$. For the monolithic CRYSP detectors $\sigma_{xyz}\sim\SI{1.7}{mm}$ (a
continuous \mbox{3-D} readout). Pixelated detectors report a fixed crystal rather than a
continuous position and are deferred.

\paragraph{Energy.}
The deposited energy is smeared by an energy-dependent resolution. The fractional FWHM
scales as $\sqrt{\SI{511}{keV}/E}$, so the absolute width is
\begin{equation}
\mathrm{FWHM}(E) = a\,\sqrt{\SI{511}{keV}\cdot E}, \qquad
\sigma_E(E) = \frac{\mathrm{FWHM}(E)}{2\sqrt{2\ln 2}},
\end{equation}
with $a$ the fractional FWHM at \SI{511}{keV}: $a=\SI{5}{\percent}$ for pure (cryogenic) CsI,
\SI{7}{\percent} for CsI(Tl), \SI{10}{\percent} for cryogenic BGO (Sec.~\ref{sec:configs}).

\paragraph{Energy window / selection.}
A hit is kept if its (smeared) energy passes the selection: a symmetric window about
\SI{511}{keV} of half-width $\sim 2\,\mathrm{FWHM}$ (tunable), and/or a minimum-energy cut.
Because the window width scales with the resolution, a sharper detector rejects more
patient scatter --- a direct lever connecting detector quality to background.

\paragraph{The clean-coincidence condition.}
For each annihilation the code looks at the blocks the two photons lit up. The clean case is
\emph{both} photons contained in one block each (no overspill) and each passing the energy
selection: the pair is recorded as a coincidence, labelled \texttt{true} if neither photon
scattered in the phantom and \texttt{scatter} if one did. Anything else fails and the event
yields nothing (an efficiency loss): a photon that missed the ring, escaped the crystal with
too little energy, or overspilled into a third block. Every detected hit is \emph{also}
written to a singles list, with its position, energy, parent event and scatter flag --- the
substrate for the randoms pass.

\paragraph{Per-crystal scintillation properties.}
The light yield, decay-time mixture, energy resolution $a$ and photodetection efficiency
(PDE) are carried on the crystal \texttt{Material} (\texttt{data/materials.json}), keyed to
the crystal name, so a run needs no separate detector cards for them. These feed both the
energy resolution and the timing model below.

% =====================================================================
\section{Timing and randoms}
\label{sec:timing}

\subsection{The timestamp: a physical first-photon model}
A gamma's timestamp is \emph{not} a free Gaussian $\sigma_t$. It is the arrival time of the
gamma's \emph{first detected scintillation photon}, a physical distribution set by the
crystal (\texttt{src/timing.jl}). A deposit of energy $E$ produces
\begin{equation}
N_{\mathrm{det}} = (\text{light yield})\cdot E \cdot \mathrm{PDE}
\end{equation}
detected photons (e.g. cryogenic CsI $\sim 10^5\,\gamma/\mathrm{MeV}\times 0.5\,\mathrm{MeV}
\times \mathrm{PDE} \sim 10^4$). Each is emitted at the scintillation mixture's effective
initial rate $r_0=\sum_k w_k/\tau_k$, and the \emph{first} of $N_{\mathrm{det}}$ such
arrivals is exponential with rate $N_{\mathrm{det}}\,r_0$, so
\begin{equation}
\text{jitter} = -\frac{\ln u}{N_{\mathrm{det}}\,r_0}, \qquad
\langle\text{jitter}\rangle = \frac{1}{N_{\mathrm{det}}\,r_0}\;\sim\;\SI{0.1}{ns}.
\end{equation}
The timing resolution thus \emph{falls out} of yield, PDE, decay time and deposited energy
--- there is no tunable $\sigma_t$. The full timestamp of a detected gamma is
\begin{equation}
t_0 = t_{\mathrm{annih}} + \mathrm{TOF} + \text{jitter}, \qquad
\mathrm{TOF} = \frac{\lVert \mathbf{hit}-\mathbf{emit}\rVert}{c},
\quad c = \SI{299.792458}{mm/ns},
\end{equation}
where $t_{\mathrm{annih}}$ is the event's annihilation time drawn from the activity curve
(\texttt{src/activity.jl}): a truncated exponential $\propto e^{-\lambda(t-t_0)}$ over the
acquisition window, with $\lambda=\ln 2/T_{1/2}$. The draw is deterministic per event number
(a splitmix64 finalizer of the seed and event id), so it is independent of how the singles
were chunked across threads and identical in every pass that reads it. The front-loading of
the activity is the origin of the randoms' front-loading.

\subsection{The three LOR categories}
The list-mode output carries a truth label per coincidence:
\begin{itemize}[nosep]
\item \texttt{true} (code 0) --- both photons reached the detector without scattering in the
phantom;
\item \texttt{scatter} (code 1) --- a same-annihilation coincidence in which at least one
photon Compton-scattered in the phantom (the patient-scatter background);
\item \texttt{random} (code 2) --- two photons from \emph{different} annihilations whose
timestamps happen to fall within the coincidence window.
\end{itemize}
For a true pair (common emission point) the timing-resolution residual is formed by
subtracting the geometric time-of-flight difference,
\begin{equation}
\Delta t = \bigl|t_{0,1}-t_{0,2}\bigr| - \mathrm{TOF}_{\mathrm{diff}}, \qquad
\mathrm{TOF}_{\mathrm{diff}} = \frac{\lVert \mathbf{hit}_1-\mathbf{emit}\rVert
- \lVert \mathbf{hit}_2-\mathbf{emit}\rVert}{c},
\end{equation}
which would be zero with no fluctuations; its spread \emph{is} the timing resolution. The
coincidence window $\tau$ is read off the trues' $|\Delta t_0|$ distribution (chosen for high
true-pair acceptance) and that same window, applied to the singles stream, lets accidental
cross-event singles leak in as randoms.

\subsection{Forming randoms}
Randoms are added in a separate pass over the singles list, with \emph{no further transport}
(\texttt{src/randoms.jl}). Each surviving single is given an absolute timestamp
($t_{\mathrm{annih}}$ of its decay $+$ its decay-relative time); the singles are sorted in
time, and any two from \emph{different} decays whose times fall within $\tau$ are emitted as a
random LOR. The implementation is an $O(n\log n)$ sort followed by a forward sliding window
that breaks as soon as the gap exceeds $\tau$ (the window holds only $\sim S\tau$ singles).
The expected rate follows
\begin{equation}
R_{\mathrm{random}} \approx 2\,\tau\,S(t)^2,
\end{equation}
with $S$ the total singles rate: front-loaded (it scales as activity squared, faster than
the trues), largest at the start of the scan, and at most a few to ten percent of the trues
for $\tau$ of a few ns. Crucially, transport --- the costly part --- runs \emph{once} to
produce the trues and the singles; randoms (and any re-realisation the downstream analysis
needs) are cheap operations on the singles list.

% =====================================================================
\section{Output and the production chain}
\label{sec:output}

\paragraph{The list-mode deliverable.}
The output is the coincidence (LOR) list: per accepted coincidence, the two \mbox{3-D} hit
positions [\si{mm}], the two energies [\si{keV}], the two timestamps [\si{ns}], the
timing-resolution residual $\Delta t$, the common emission point, and the truth label
(\texttt{true}/\texttt{scatter}/\texttt{random}). A companion \texttt{\_meta} record carries
the provenance: detector config, geometry, energy and time windows, scenario name and seed.
This list is what the deferred downstream analysis reads, and the only thing it reads.

\paragraph{The chain.}
The production chain (\texttt{scripts/}) runs in two stages. First a multi-threaded transport
(\texttt{simulate\_source\_mt.jl}, via the allocation-free single-photon navigator) writes a
\emph{singles} stack (\texttt{singles.\{csv,h5\}}). Then a streaming selection
(\texttt{build\_coincidences\_from\_singles.jl}) reads the singles and writes the LOR file
(\texttt{lors\_\{truth,det\}.h5}); the randoms pass runs over the same singles. The selection
is a \emph{streaming}, event-ordered pass with $O(1)$ memory --- it never loads the whole
file --- which is why it is written in Julia rather than Python: it must scale to the large
stacks a full run produces. HDF5 stores quantized \texttt{Int16} columns
(\SI{0.1}{mm}/\SI{0.1}{keV}, lossless at detector resolution), chunked and compressed
($\sim 6\times$ smaller than float CSV, with typed partial reads) for the write-once,
read-many singles list. CSV is kept for development and inspection. Each run is driven by a
TOML config (\texttt{runs/<tag>.toml}, sections
\texttt{[geometry] [source] [transport] [detector] [output]}) that is both the parameter
source of truth and the run's provenance; its tag names the per-run output directory.

\paragraph{Division of labour.}
Julia carries the photon transport, the geometry, and the streaming coincidence/randoms
selection (memory and throughput on large stacks). Python reads the scenario and makes the
control and analysis plots. CSV and HDF5 are the interchange formats.

% =====================================================================
\section{Detector configurations under study}
\label{sec:configs}

The detectors of present interest are \emph{monolithic} crystals with continuous \mbox{3-D}
readout (CRYSP; numbers from Soleti et al.\ 2024):
\begin{center}
\begin{tabular}{lccc}
\hline
Crystal & $a$ (FWHM @ \SI{511}{keV}) & light yield [$\gamma$/MeV] & decay [\si{ns}] \\
\hline
pure CsI (cryogenic) & \SI{5}{\percent}  & $1.0\times10^5$ & 1000 (single) \\
CsI(Tl)              & \SI{7}{\percent}  & --- & --- \\
BGO (cryogenic)      & \SI{10}{\percent} & $2.9\times10^4$ & 747 / 2253 (0.34 / 0.66) \\
\hline
\end{tabular}
\end{center}
All share a continuous position/DOI resolution $\sigma_{xyz}\sim\SI{1.7}{mm}$. The PDE is
currently a \num{0.45} placeholder for both CsI and BGO and should be refined per crystal
(it depends on the photodetector efficiency at the crystal's emission wavelength --- CsI in
the UV, BGO near \SI{480}{nm}). Pixelated detectors (LYSO, standard BGO, $a\sim$
\SIrange{15}{20}{\percent}) report a fixed crystal rather than a continuous position and come
later. Because every detector reads the identical scenario through the identical transport,
the differences in their coincidence lists are due to the crystal and readout alone --- which
is precisely what makes the project a controlled comparison bench.

% =====================================================================
\section*{Reference material}
\addcontentsline{toc}{section}{Reference material}
\begin{itemize}[nosep]
\item Method guide (companion to this overview): \texttt{docs/pet\_simulation.tex}.
\item Ray--solid navigation math: \texttt{docs/navigation.tex}.
\item The randoms / timing plan: \texttt{dev/randoms\_plan.md}.
\item CRYSP detector numbers: Soleti et al.\ 2024.
\item Upstream source method: \texttt{ptcryspg4/docs/simulate\_pt\_pet.tex}.
\end{itemize}

\end{document}
